
\documentclass[12pt]{article}
\usepackage{multirow}
\usepackage[pdftex]{graphicx}
\usepackage{float}
\newenvironment{my_itemize}
{\begin{itemize}
  \setlength{\itemsep}{0cm}
  \setlength{\parskip}{0cm}}
{\end{itemize}}

\begin{document}

\begin{titlepage}
\vspace{-1.5cm}
\begin{center}
\includegraphics[width=2cm]{Logo_Alexandria_University.jpg}\\
\vspace{1cm}
\textbf{\large ALEXANDRIA UNIVERSITY} \\
\textbf{FACULTY OF ENGINEERING} \\
{\small  COMPUTER AND SYSTEMS ENGINEERING DEPARTMENT}

\vspace{2.5cm}
\textbf{\LARGE Taxonomy for Email Classification and Summarization Techniques}\\
\vspace{1cm}
{ Ahmed El-Sharkasy, Ahmed Kotb, Amr Nabil, Mohammad Kotb, Moustafa Mahmoud }
\end{center}

\vspace{1ex}
\textbf{Supervisors:} Prof. Dr. Mohamed Abou-gabal, Dr. Mustafa Elnainany
\end{titlepage}

\newpage
\tableofcontents
\newpage

\section{Abstract}
In this document we present a survey and taxonomy on recent research topics 
related to email classification and summarization. This document summarizes 
and organizes recent research results in the novel way that integrates and 
adds understanding to work in the field of email classification and 
summarization. It emphasizes the classification of the existing literature, 
developing a perspective on the area, and evaluating different trends.

\paragraph{Keywords}
Email, Classification, Summarization, Machine Learning.

\section{Introduction}
Email has been an efficient and popular communication mechanism as the 
number of Internet users increases. Therefore, email management has become 
an important and growing problem for individuals and organizations because 
it is prone to misuse. One of the problems that are most paramount is disordered 
email message, congested and unstructured emails in mail boxes. It may be very 
hard to find archived email message, search for previous emails with specified 
contents or features when the mails are not well structured and organized.

Many machine learning approaches have been applied in this field, the most 
State-of-the-Art algorithms in email classification include: support vector 
machines, neural network, naïve bayes classifiers and entropy-based approach. 

Email summarization is another important and challenging problem. We can think 
of automatic summarization as a type of information compression. To achieve such 
compression, better modelling and understanding of document structures and internal 
relations is required. 

In this document we present a survey and taxonomy on recent research topics 
related to email classification and summarization.

\section{Email Classification Taxonomy}
The following table classifies some recent research papers in the field of email 
classification according to the different learning algorithms used in different papers


\begin{tabular}{|p{2cm}|p{2cm}|p{2cm}|p{2cm}|p{2cm}|p{2cm}|}
\hline
\multicolumn{6}{|c|}{Learning Algorithm} \\
\hline
SVM & Naïve Bayes & Neural Networks & Max. Entropy / Winnow & Nnge / Hoeffing Trees & Graph Mining \\ \hline
An Innovative Analyser for email classification Based on Grey List Analysis &
Email Classification with Co-training &
Email Classification: Solution with Back Propagation Technique & 
Automatic Categorization of Emails into Folders &
Using GNUsmail to compare Data Stream Mining Methods for On-line Email Classification &
A graph Based Approach for Multi-Folder Email Classification \\ \hline

Email Classification with Co-training &
Automatic Categorization of Emails into Folders &
Email Classification Using Semantic Feature Space & 
&
&
 \\ \hline

Automatic Categorization of Emails into Folders &
&
& 
&
&

 \\
\hline
\end{tabular}

\section{Email Summarization Taxonomy}

\section{Papers Summary}
\subsection{Email Classification}

\subsubsection{Automatic Categorization of Email into Folders \cite{RON04}}

\paragraph{Year} 2004
\paragraph{Citations} 112
\paragraph{Introduction}
\begin{my_itemize}
  \item Users get alot of emails this days, not just spam but a large number of 
	legitmate emails also that they need to process in a short time.
  \item The paper shows the results of an extensive benchmark on two large corpora 
	(enron,sri) of 4 classification algorithms.
  \item The paper shows an enhancement to the exponential gradient method (winnow).
\end{my_itemize}

\paragraph{Related Work}
\begin{my_itemize}
  \item Clark and Niblet 1989: proposed a rule inductive algorithm CN2 and 
	showed that it can outperform KNN.
  \item Cohen 1996: proposed the RIPPER classifier and showed that that it 
	can outperfrom an tfidf classifier.
  \item Provost 1999: showed that Naive bayes can outperform RIPPER.
  \item Remmie 2000: achived a very high accuracy by classifying mails to 
	3 predefined folders .
  \item Kiritchenko and Malwin 2001: showed that SVM can outperfom Naive Bayes.
\end{my_itemize}


\paragraph{Algorithms Benchmarked}
\begin{my_itemize}
  \item Maximum Entropy.
  \item Naive Bayes.
  \item SVM.
  \item Winnow (enhanced version).
\end{my_itemize}

\paragraph{Challenges in mail classification}
\begin{my_itemize}
  \item Email users often create folders and let it fall out of use 
	(small number of training data per folder).
  \item Folders don’t necessarily correspond to simple semantic topics 
	(unfinished todos, project groups, certain recipient).
  \item Differ drastically from one user to another.
  \item Email arrives in a stream over time which causes more difficulties, 
	for example the topic of main folder can drift over time.
\end{my_itemize}


\paragraph{Data set pre-processing}
\begin{my_itemize}
    \item Removing non topical folders (Inbox, sent, trash, ...etc).
    \item Removing small folders (folders that has a small number of emails).
\end{my_itemize}

\paragraph{Training/test set splits}
\begin{my_itemize}
    \item The paper shows a new way to split training data into training set 
	  and test sit, the new method takes time factor into considerations.
    \item It works as follows:
    \begin{my_itemize}
        \item sorting emails by time;
        \item train the classifier for the first N emails;
        \item test it on the following N emails;
        \item train the classifier for the first 2N emails;
        \item then test it for the following N emails;
        \item and so on.
    \end{my_itemize}
\end{my_itemize}

\paragraph{Features Extraction}
traditional bag of words representation.

\paragraph{Datasets}
    \begin{my_itemize}
    \item Enron: http://www.cs.cmu.edu/~enron/
    \begin{my_itemize}
        \item 150 users with more than 500,000 Emails;
        \item applied to the following 7 employees folder only 
	      (the largest 7 folders : beck-s,farmer-d, kaminski-v, 
	      kitchen-l, lokay-m, sanders-r and williams-w3);
        \item removed the non topical folders like ``all documents'', 
	      ``calendar'', ``contacts'', ``deleted items'', ``discussion threads'', 
	      ``inbox'', ``notes inbox'', ``sent'', ``sent items'' and ``sent mail'';
        \item flatten all the folder hierarchies;
        \item removed folders with less than 3 messages;
        \item removed the X-Folder field from email messages. (The X-Folder 
	      field contains the class label).
    \end{my_itemize}
    \item SRI : http://www.ai.sri.com/project/CALO
    \begin{my_itemize}
        \item applied to the following 7 folders only: acheyer, bmark, disrael, 
	      mgervasio, mgondek, rperrault, and vchaudri;
        \item removed the non topical folders (inbox, draft, sent, trash);
        \item flatten all the folder hierarchies;
        \item removed folders with less than 3 messages.
    \end{my_itemize}
\end{my_itemize}

\paragraph{Critique}
\begin{my_itemize}
    \item They didn’t use Stemming in their preprocessing to the dataset.
    \item Not including precision , recall and f1 score for accuracy measures.
\end{my_itemize}

\paragraph{Conclusion}
\begin{my_itemize}
    \item Naive Bayes is inferior to other algorithms.
    \item SVM achieved the highest accuracies in most of the tests.
\end{my_itemize}

\paragraph{Future Work}
\begin{my_itemize}
    \item Different sections of each email can be treated differently. 
	  For example, the system could create distinct features for words appearing 
	  in the header, body, signature, attachments, ...etc.
    \item Named entities may be highly relevant features. It would be desirable to 
	  incorporate a named entity extractor (such as MinorThird3, see, e.g., 
	  Cohen and Sarawagi (2004)) into the foldering system.
\end{my_itemize}

%==============================================================================

\subsubsection{Email Classifications For Contact Centers \cite{ANI03}}
\paragraph{Year} 2003
\paragraph{Citations} 14
\paragraph{Main Topic}
\begin{my_itemize}
    \item Proposing an automatic system to classify mail message for contact centers.
    \item Mails are categorized into 2 classes:
    \begin{my_itemize}
        \item single messages: messages that don’t require a response;
        \item root messages: messages that require immediate response;
        \item root messages can be sub divided into 3 classes:
        \begin{my_itemize}
            \item root: the start of the communication (contains a problem or a question);
            \item inner: communication on a certain problem;
            \item leaf: marks the end of this interaction (eg. the problem was solved).
        \end{my_itemize}
    \end{my_itemize}
\end{my_itemize}

\paragraph{Tools used}
\begin{my_itemize}
    \item Rainbow: an implementation for naive bayess algorithm.
    \item SVMlight: an implementation for SVM algorithm.
    \item WordNet: used for parts of speach taging.
    \item Ltchunk: used to identify noun phrases and count number of sentences in email.
\end{my_itemize}

\paragraph{Dataset}
\begin{my_itemize}
    \item Pine-info discussion list web archive
    \begin{my_itemize}
        \item http://www.washington.edu/pine/pine-info.
    \end{my_itemize}
\end{my_itemize}

\paragraph{Pre-processing}
\begin{my_itemize}
    \item Removing reply blocks (blocks from previous emails in the current mail).
    \item Removing signature blocks.
\end{my_itemize}

\paragraph{Features (for SVM algorithm)}
\begin{my_itemize}
    \item Non-infected words
    \begin{my_itemize}
        \item nouns, verbs, adjective, adverb;
        \item using WordNet;
    \end{my_itemize}
    \item Noun phrases
    \begin{my_itemize}
        \item using Ltchunk;
    \end{my_itemize}
    \item Verb phrases.
    \item Punctuation letters count.
    \item Length of email (number of sentences)
    \begin{my_itemize}
        \item using Ltchunk;
    \end{my_itemize}
    \item Dictionary
    \begin{my_itemize}
        \item 2 dictionaries were made one for the most common words in single 
	      messages and the other for the most common words in root message.
    \end{my_itemize}
\end{my_itemize}

\paragraph{Conclusion}
\begin{my_itemize}
    \item High accuracy was achieved on root vs leaf (92\%) , root vs inner (87\%) and root vs single(79\%).
\end{my_itemize}


%=================================================================================================

\subsubsection{Using GNUsmail to Compare Data Stream Mining Methods for On-line Email \cite{JOSE11}}
\paragraph{Year} 2011

\paragraph{Citations} 0

\paragraph{Main Topic}
\begin{my_itemize}
    \item Introducing GNUsmail, an open source framework used for mail 
	  classification, focusing on online incremental learning.
    \item Proposing new techniques for testing other than holdout and 
	  cross-validation like prequential measure.
\end{my_itemize}

\paragraph{Evaluation methods}
\begin{my_itemize}
    \item Prequential measure.
    \item Sliding and fading windows.
    \item McNemar test.
\end{my_itemize}

\paragraph{Dataset}
\begin{my_itemize}
    \item Enron.
    \item A layer was added to feed the learning algorithm the new emails one 
	  by one, simulating new incoming emails.
\end{my_itemize}

\paragraph{Algorithms}
\begin{my_itemize}
    \item OzaBag over NNge, using DDM for concept drift detection.
    \item NNge.
    \item Hoeffding Trees.
    \item Majority class.
\end{my_itemize}

\paragraph{Tools}
\begin{my_itemize}
    \item GNUsmail: http://code.google.com/p/gnusmail/.
\end{my_itemize}

\paragraph{Result}
\begin{my_itemize}
  \item Improved GNUsmail by incorporating new different methods to evaluate 
	data stream mining algorithms in the domain of email classification.
\end{my_itemize}

\paragraph{Future Work}
\begin{my_itemize}
  \item Current online learning algorithm implementations have an important 
	limitation that affects the learning process: learning attributes have 
	to be fixed before beginning the induction of the algorithm. They need 
	to know all the attributes, values and classes before the learning itself, 
	since it is not possible to start using a new attribute in the middle of the 
	lifetime of a learning model. Future methods should support online addition of
	new features.
\end{my_itemize}

%=================================================================================================

\subsubsection{E-Classifier: A Bi-Lingual Email Classification System \cite{NOUF08}}
\paragraph{Year} 2008
\paragraph{Citations} 0

\paragraph{Problem}
\begin{my_itemize}
    \item Classifying Arabic and English emails.
    \item Implementing an outlook add-in ``e-classifier''.
\end{my_itemize}

\paragraph{Related Work}
\begin{my_itemize}
    \item English Email Classifiers
    \begin{my_itemize}
        \item PopFile
        \begin{my_itemize}
            \item http://popfile.sourceforge.net.
            \item Uses naive bayes algorithm only.
        \end{my_itemize}
        \item SpamBayes
        \begin{my_itemize}
            \item http://spambayes.sourceforge.net
            \item Binary Classifier (Spam or not).
        \end{my_itemize}
    \end{my_itemize}
    \item Arabic Email Classifiers
    \begin{my_itemize}
    \item There are no Email classification work on arabic language, the
	  related work are on arabic documents not emails El-Kourdiet.
    \end{my_itemize}
\end{my_itemize}

\paragraph{Dataset}
\begin{my_itemize}
    \item English: enron dataset.
    \item Arabic
    \begin{my_itemize}
        \item Translated documents that have been converted to emails.
        \item Documents obtained from http://www.comp.leeds.ac.uk/eric/latifa/research.html.
    \end{my_itemize}
\end{my_itemize}

\paragraph{pre-processing}
\begin{my_itemize}
    \item English
    \begin{my_itemize}
        \item Removing stop words.
        \item Removing punctuation marks.
        \item Converting all the letters to lowercase.
        \item Porter stemmer.
    \end{my_itemize}
    \item Arabic
    \begin{my_itemize}
        \item No root extraction technique was used due to the lack of non commercial product.
    \end{my_itemize}
\end{my_itemize}

\paragraph{Results}
\begin{my_itemize}
    \item 85\% of English emails were classified correctly.
    \item 60\% of Arabic emails were classified correctly.
\end{my_itemize}

\paragraph{Critique}
\begin{my_itemize}
    \item Used only overall accuracy measure which might not good indicator in case of skewed data.
\end{my_itemize}

%=================================================================================================

\subsubsection{An Object Oriented Email Clustering Model Using Weighted 
	      Similarities between Emails Attributes \cite{NARESH10}}
\paragraph{Year} 2010
\paragraph{Citations}6
\paragraph{Description}
\begin{my_itemize}
    \item Proposing a new Object Oriented Email Clustering Model to categorize 
	  mail message into groups.
\end{my_itemize}

\paragraph{Algorithms}
\begin{my_itemize}
    \item K-means clustering algorithms.
    \item Text similarity techniques.
    \begin{my_itemize}
        \item cosine Similarity.
        \item Dice Similarity.
        \item Blue Similarity.
        \item TF-IDF Similairty (Term Frequency - Inverse Domain Frequency).
        \item Jaccard Similairty.
    \end{my_itemize}
\end{my_itemize}

\paragraph{Datasets}
\begin{my_itemize}
    \item Enron dataset.
    \item Inbox folder of base-e user mail box.
\end{my_itemize}

\paragraph{Dataset pre-processing}
\begin{my_itemize}
    \item Stemming.
    \item Parsing
    \begin{my_itemize}
        \item To extract email attribuites (subject, body, ...etc).
    \end{my_itemize}
    \item Storing in an object oriented representation.
\end{my_itemize}

\paragraph{Tools/programming languages used}
\begin{my_itemize}
    \item Java.
    \item Simmetric: used to calculate text similarities.
    \item Weka (Waikato Environment for Knowledge Analysis): used for stemming of emails.
\end{my_itemize}

\paragraph{Future work}
\begin{my_itemize}
    \item Thread summarization.
    \item Automatic email answering.
\end{my_itemize}

\paragraph{Conclusion}
\begin{my_itemize}
    \item Email can be represented as an object with attribute like subject, body, ...etc.
    \item Clustering of emails can be implemented in an object oriented way.
\end{my_itemize}


%======================================================================================

\subsubsection{Content Based Email Classification System by applying Conceptual Maps \cite{BASKARAN09}}

\paragraph{Year} 2009
\paragraph{Citations} 0

\paragraph{Main Topic}
\begin{my_itemize}
    \item Proposing a Knowledge based System (KBS) to classify messages into folders.
    \item Using lexicon and conceptual graphs.
\end{my_itemize}

\paragraph{Major steps of processing on subject and body fields}
\begin{my_itemize}
    \item Word splitting.
    \item Word normalization (stemming).
    \item Detect abbreviation.
    \item Removing stop words.
    \item Word indexing.
    \item Identify noun-phrases by NLP techniques.
    \item Conversion of phrases into concepts.
\end{my_itemize}

\paragraph{Related Work}
\begin{my_itemize}
    \item C-Evolove.
    \item Titus.
\end{my_itemize}

%=============================================================================================
\subsubsection{A new approach to Email classification using Concept Vector Space Model \cite{CHAO08}} 

\paragraph{Year} 2008
\paragraph{Citations} 3
\paragraph{Algorithms}
\begin{my_itemize}
  \item Used a classification algorithms based on pre-processing steps in the 
	training phase to produce vector that identify the category or the new email.
  \item Based on WordNet, for describing a text Email by establishing concept vector
	space model, we can firstly extract the high-level information on categories
	during training process by replacing terms with synonymy sets in WordNet and
	considering hypernymy-hyponymy relation between synonymy sets.
  \item Used TF * IWF * IWF method to revise the weight of the concept vector.
\end{my_itemize}

\paragraph{Dataset}
\begin{my_itemize} \setlength{\itemsep}{0cm}%
  \setlength{\parskip}{0cm}%
  \item Used documents of 20 news group (standard document set).
  \item Put these documents in 20 directory as 20 category, each category contains at least 1,000 article.
  \item Set of these articles are selected to be used as training set, another set as a test set.
\end{my_itemize}



\paragraph{Results}
\begin{my_itemize}
  \item Made two experiments on different conditions, comparing the concept VSM method with a traditional VSM method:
  \begin{my_itemize}
    \item experiment 1:
    \begin{my_itemize}
      \item selected 3 categories from the dataset, and chosen 300 email at random 
	    from each category as training set and 100 email from each category as 
	    test set;
      \item observed that the F1-meausre of the concept VSM is always better than 
	    traditional VSM by at least factor of 0.1 (for more details check 
	    Tables 1,2 in the paper) with F1-measure for concept VSM in the 3 
	    datasets 0.84, 0.90, 0.93 respectively.
    \end{my_itemize}
    \item experiment 2:
    \begin{my_itemize}
      \item used the same categories in experiment one, but repeated experiment 
	    one but with different training set size, starting from 30 email;
      \item observed that Concept VSM is always better than traditional VSM;
      \item accuracy starts from 0.4 at 30 email training set for all categories 
	    and increases till it reach 0.9 for training set size as in 
	    experiment 1 (900 emails for all categories);
      \item this means that Concept VSM is working fine with small training set 
	    size but it is better if it is increased;
      \item for more details check Figure 2 in the paper.
    \end{my_itemize}
  \end{my_itemize}
\end{my_itemize}

\paragraph{Future work}
Use concept VSM to do level classification.

%=============================================================================================
\subsubsection{Ontology based classification and categorization of email \cite{BALAKUMAR08}}

\paragraph{Year} 2008
\paragraph{Citations} 4

\paragraph{Problem}
\begin{my_itemize}
  \item Making a user defined and user controllable spam filter to detect spam emails, 
	the paper uses ontology for understanding the content of the email and Bayesian
	approach for making the classification.
  \item Categorizing mails based on their content.
  \item The complete process: classifying mails as hams or spams and further classification of ham emails to folders.
\end{my_itemize}

\paragraph{Algorithms}
\begin{my_itemize}
  \item Content based filtering: uses keywords in the mail for classification.
  \item Statistical based filtering: Assigns probability or score to each keyword 
	and uses the overall probability or score to classify the new mail.
  \item Machine learning approach for filtering: Ontology is used as one of the 
	learning tools for email classification.
\end{my_itemize}

\paragraph{Results}
\begin{my_itemize}
  \item 98\% of the emails has been classified successfully to ham and spam.
  \item 95\% of the ham has been successfully categorized into folders.
\end{my_itemize}

\paragraph{Conclusion}
\begin{my_itemize}
  \item User defined spam filter has better results than general spam filters for all user.
\end{my_itemize}

%=============================================================================================
\subsubsection{Enterprise Email Classification Based on Social Network Features \cite{MIN11}}

\paragraph{Year} 2011
\paragraph{Citations} 0

\paragraph{Problem}
Managing the email services in Enterprises, so that business emails have priority 
over personal emails by classifying emails into official and private. The 
classification is made based on social features not on the email content for 
protecting the privacy of users' emails by building a social network analysis 
graph representing the senders and recipients as vertixes and the sending events as edges.

\paragraph{Algorithms}
\begin{my_itemize}
  \item Support vector machine (SVM).
  \item WEKA.
\end{my_itemize}


\paragraph{Results}
\begin{my_itemize}
  \item F-measure = 0.9
\end{my_itemize}

\paragraph{Related Work}
\begin{my_itemize}
  \item SNARF http://research.microsoft.com/en-us/projects/snarf/
\end{my_itemize}

\paragraph{Future work}
\begin{my_itemize}
  \item Combining some state-of–the-art email prioritization algorithms with the 
	proposed method to balance the loading of email server.
\end{my_itemize}


%=============================================================================================
\subsubsection{Email Categorization Using Multi-Stage Classification Technique \cite{MD07}}

\paragraph{Year} 2007
\paragraph{Citations} 5

\paragraph{Problem}
\begin{my_itemize}
  \item Email classification (spam) using a multi-stage classification technique 
	collecting all the mails which is not TP or TN in a different mailbox for
	the user to give feedback about them. The classification of emails is done 
	in multi stages where in each stage a new classifier is added to filter output.
\end{my_itemize}

\paragraph{Algorithms}
\begin{my_itemize}
  \item SVM.
  \item Naive Bayes.
  \item Boosting Algorithms.
\end{my_itemize}

\paragraph{Results}
\begin{my_itemize}
  \item Average FP is 0 and average FN is lower than that results from using any algorithm individually.
  \item Accuracy : 97.05\%.
\end{my_itemize}

\paragraph{Data sets}
\begin{my_itemize}
  \item PUA.
\end{my_itemize}

\paragraph{Future work}
\begin{my_itemize}
  \item Analyse cost in complexity and speed.
\end{my_itemize}


%=============================================================================================
\subsubsection{Automatically tagging email by leveraging other users folders \cite{YEHUDA11}}

\paragraph{Year} 2011
\paragraph{Citations} 0

\paragraph{Problem}
\begin{my_itemize}
  \item Automatically associating semantic tags to emails other than creating folders. 
	Beside providing a way to tag emails automatically, they started with predefined 
	set of tags taken from a study on the folders and labels yahoo users are generating. 
	The proposed technique took into consideration the performance and scalability. 
	The technique learned how to tag by taking into account the habit of many users for 
	making folders simultaneously.
\end{my_itemize}

\paragraph{Algorithms}
\begin{my_itemize}
  \item K-means.
  \item Naive bayes.
  \item Other proposed algorithms.
\end{my_itemize}

\paragraph{Data sets}
\begin{my_itemize}
  \item Emails from yahoo mail users (200 million emails).
\end{my_itemize}


\paragraph{Conclusion}
\begin{my_itemize}
  \item The paper presented a classification system for tagging emails suitable for a very 
	large scale system up to millions of emails and reached a performance of 2 ms or 
	less for classifying an email and with acceptable accuracy.
\end{my_itemize}

\paragraph{Future work}
\begin{my_itemize}
  \item Increasing the features extracted from the emails to include To: and Cc: fields,
	the length of a message, the number and names of file attachments,style (html/plain) signals, 
	and more sophisticated subject tokenization techniques.
\end{my_itemize}



%=============================================================================================
\subsubsection{An Email Classification Model Based on Rough Set Theory \cite{WENQING05}}

\paragraph{Year} 2005
\paragraph{Citations} 19

\paragraph{Problem}
\begin{my_itemize}
  \item Reducing the error rate of classifying non-spam emails into spam by classifying 
	the incoming emails into 3 categories instead of 2: spam, non-spam and suspicious
	using an algorithm based on rough set theory.
\end{my_itemize}

\paragraph{Related work}
\begin{my_itemize}
  \item Ripper Algorithm.
  \item Genetic Document Classifier.
  \item Smokey.
  \item Bayesian Junk Email Filter.
  \item Max. Entropy Model.
\end{my_itemize}

\paragraph{Data sets}
\begin{my_itemize}
  \item http://www.ics.uci.edu/mlearn/MLRepository.html
\end{my_itemize}

\paragraph{Results}
\begin{my_itemize}
  \item Accuracy reached 97\%.
\end{my_itemize}

\paragraph{Conclusion}
\begin{my_itemize}
  \item Rough set based model can reduce the error rate that classifies a non-spam email to spam.
\end{my_itemize}

%=============================================================================================
\subsection{Email Summarization}
\subsubsection{Detection of question-answer pairs in email conversations \cite{LOKESH04}}

\paragraph{Year} 2004
\paragraph{Citations} 41

\paragraph{Problem}
\begin{my_itemize}
  \item The sentence extraction summarization method can’t be applied in all 
	types of documents.
  \item Using it in summarizing email threads is not efficient, as it is a very
	special type of documents, as sentences and words are written relative 
	to previous emails, so using sentence extraction will not be useful in this
	case.
  \item This paper is trying to solve this problem by extracting pairs of 
	questions and answers to summarize email threads.
\end{my_itemize}

\paragraph{Conclusion}
\begin{my_itemize}
  \item Good approach to extract question-answer pairs in the email conversation 
	in case of interrogative questions.
  \item Declarative and rhetorical questions can’t be detected such as 
	``Please let  me know ...'', ``I was wondering if ...'', 
	``If you could ..., that would be great''.
  \item Future work is to investigate these types of questions.
\end{my_itemize}



%=============================================================================================

\subsubsection{Using Question-Answer Pairs in Extractive Summarization of Email Conversations \cite{KATHLEEN07}}

\paragraph{Year} 2007
\paragraph{Citations} 12

\paragraph{Problem}
\begin{my_itemize}
  \item After 3 years from the previous paper, they thought for a new approach 
	to make a hybrid solution by extractive summarization of email threads 
	with automatically detected QA pairs.
  \item This approach is better than extracting QA pairs only, as due to some 
	statistics they made on their dataset that:
  \begin{my_itemize}
    \item 20\% of emails are question-answer exchange;
    \item 40\% of all email threads involve question-answer exchange of some form.
  \end{my_itemize}
  \item Sentence extraction may be very useful if augmented with email specific 
	features as dialogic structure.
\end{my_itemize}

\paragraph{Algorithms}
\begin{my_itemize}
  \item Extractive summarization:
  \begin{my_itemize}
    \item represent each sentence in the SEQA threadset with a feature vector 
	  along with its binary classification, which represents wether or not a
	  sentence should be in a summary;
    \item features used are length, position in the document, TF-IDF scores, ...etc.
  \end{my_itemize}
  \item QA Pair detection:
  \begin{my_itemize}
    \item train a classifer on QA detection on the data corpus.
  \end{my_itemize}
  \item Integrating QA Pairs with Extractive summarization:
  \begin{my_itemize}
    \item 3 different approaches:
    \begin{my_itemize}
      \item SE+A: a sentence figures as an answer to a question asked earlier 
	    in the thread as an additional feature in our machine learning-based 
	    extractive summarization approach;
      \item SE+QA: to add automatically detected answers to questions in 
	    extractive summaries and add detected questions to answers in extractive 
	    summaries not in the summaries;
      \item QA+SE: start with automatically detected question-answer pair sentences 
	    which are then augmented with extractive sentences that do not 
	    appear already in the question-answer pair sentences.
    \end{my_itemize}
  \end{my_itemize}
\end{my_itemize}


\paragraph{Data sets}
\begin{my_itemize}
  \item Corpus contains 300 email thread, each thread contains on average 3.25 email message.
  \item Data set was prepared manually concerning these points:
  \begin{my_itemize}
    \item write summaries of email threads of the corpus;
    \item highlight and link QA pairs in the email thread
    \begin{my_itemize}
      \item Highlight only the questions that seek information (wether it is 
	    interogative or declarative questions, with or without question mark, 
	    but not rhetorical questions).
      \item Link question with its answer if it was found in the same thread.
    \end{my_itemize}
  \end{my_itemize}
  \item SEQA threadset is set of email threads containing QA pairs identified manually of size 44 email thread.
\end{my_itemize}


%=============================================================================================

\subsubsection{Summarizing email conversations with clue words \cite{GIUSEPPE07}}

\paragraph{Year} 2007
\paragraph{Citations} 48

\paragraph{Problem} 
Proposing a new framework to summarize emails by capturing email conversations 
and giving weights to sentences. Their algorithm allows the user to specify the size of the summary

\paragraph{Related Work}
\begin{my_itemize}
  \item Multi-Document summarization method.
  \item Ripper classifier.
  \item And more.
\end{my_itemize}

\paragraph{Algorithms}
\begin{my_itemize}
  \item Clue word summarizer (CWS).
  \item Porter’s stemming algorithm.
  \item MEAD Summarizer.
\end{my_itemize}

\paragraph{Dataset} Enron Dataset.

\paragraph{Conclusion}
This papers introduces the fragment quotation graph which represents the 
conversation structure of the emails.This graph includes hidden emails and can 
represent the conversation in more details than a simple threading structure. 
Based on the fragment quotation graph, a new summarization approach CWS has 
been developed to select important sentences from an email conversation

\paragraph{Future Work}
Improving the fragment quotation graph generation with more sophisticated linguistic
analysis and also evaluating the algorithm with different data sets


%=============================================================================================

\section{Results}\label{results}
In this section we describe the results.

\section{Conclusions}\label{conclusions}
We worked hard, and achieved very little.


\begin{thebibliography}{99}
\bibitem{RON04}
  Ron Bekkerman,
  Andrew McCallum,
  Gary Huang,
  \emph{Automatic Categorization of Email into Folders: Benchmark Experiments on Enron and SRI Corpora},
  2004.

\bibitem{ANI03}
  Ani Nenkova,
  Amit Bagga,
  \emph{Email Classification for Contact Centers},
  2003.

\bibitem{JOSE11}
  Jose M. Carmona-Cejudo,
  Manuel Baena-Garcia,
  Jose del Campo-Avila,
  Rafael Morales-Bueno,
  Joao Gama,
  Albert Bifet,
  \emph{Using GNUsmail to Compare Data Stream Mining Methods for On-line Email Classification},
  2011.

\bibitem{NOUF08}
  Nouf Al Fe'ar,
  Einas Al Turki,
  Asma Al Zaid,
  Mashael Al Duwais,
  Mona Al Sheddi,
  Nora Al khamees,
  Nouf Al Drees,
  \emph{E-Classifier: A Bi-Lingual Email Classification System},
  2008.

\bibitem{NARESH10}
  Naresh Kumar Nagwani,
  Ashok Bhansali,
  \emph{An Object Oriented Email Clustering Model Using Weighted Similarities between Emails Attributes},
  2010.


\bibitem{BASKARAN09}
  S. Baskaran,
  \emph{Content Based Email Classification System by applying Conceptual Maps},
  2009.

\bibitem{CHAO08}
  Chao Zeng,
  Zhao Lu,
  Junzhong Gu,
  \emph{A new approach to Email classification using Concept Vector Space Model},
  2009.

\bibitem{BALAKUMAR08}
  M.Balakumar,
  V.Vaidehi,
  \emph{Ontology based classification and categorization of email},
  2008.

\bibitem{MIN11}
  Min-Feng Wang,
  Sie-Long Jheng,
  Meng-Feng Tsai,
  Cheng-Hsien Tang,
  \emph{Enterprise Email Classification Based on Social Network Features},
  2011.

\bibitem{MD07}
  Md Rafiqul Islam,
  Wanlei Zhou,
  \emph{Email Categorization Using Multi-Stage Classification Technique},
  2007.

\bibitem{YEHUDA11}
  Yehuda Koren,
  Edo Liberty,
  Yoelle Maarek,
  Roman Sandler,
  \emph{Automatically Tagging Email by Leveraging Other Users' Folders},
  2011.

\bibitem{WENQING05}
  Wenqing Zhao,
  Zili Zhang,
  \emph{An Email Classification Model Based on Rough Set Theory},
  2005.

\bibitem{LOKESH04}
  Lokesh Shrestha,
  Kathleen McKeown,
  \emph{Detection of question-answer pairs in email conversations},
  2004.

\bibitem{KATHLEEN07}
  Kathleen McKeown,
  Lokesh Shrestha,
  Owen Rambow,
  \emph{Using Question-Answer Pairs in Extractive Summarization of Email Conversations},
  2007.

\bibitem{GIUSEPPE07}
  Giuseppe Carenini,
  Raymond T. Ng,
  Xiaodong Zhou,
  \emph{Summarizing Email Conversations with Clue Words},
  2007.
\end{thebibliography}


\end{document}
